\chapter{DESCONEXIÓN EMOCIONAL EN LAS REDES SOCIALES}
El presente capítulo responde al primer objetivo específico del proyecto: analizar el nivel de desconexión emocional que experimentan los jóvenes al interactuar con plataformas digitales, identificando factores que limitan la generación de vínculos significativos. Para ello, en el presente trabajo se entiende por \textit{desconexión emocional} la percepción de que las interacciones digitales no satisfacen la necesidad de conexión genuina, manifestada en sentimientos de soledad y en vínculos de bajo significado emocional pese a la existencia de numerosos contactos. Esta definición operativa se apoya en la caracterización de la soledad como la percepción de que las propias necesidades sociales no están siendo satisfechas por la cantidad o, especialmente, por la calidad de las relaciones (Hawkley \& Cacioppo, 2010), y en la experiencia, descrita por los propios usuarios de redes sociales, de sentirse rodeados de personas y, al mismo tiempo, emocional o socialmente aislados (Parry et al., 2026).

\section{Uso de redes sociales y soledad en jóvenes}
La relación entre el uso de redes sociales y la soledad ha sido objeto de numerosas investigaciones en los últimos años, con resultados que durante mucho tiempo resultaron contradictorios. Sin embargo, la evidencia más reciente permite precisar en qué condiciones esta relación se manifiesta.
Feng et al. (2026) realizaron un meta-análisis que sintetizó 59 estudios con 67 muestras independientes y un total de 81,344 participantes adolescentes y jóvenes adultos (edad media de 17.6 años), publicados entre el año 2000 y mayo de 2025. Los autores reportan que ``el uso problemático de las redes sociales mostró una asociación moderada y consistente con una mayor soledad ($r = 0.29$, IC 95\,\% [0.22--0.35], $p < 0.001$), mientras que el tiempo dedicado a las redes sociales mostró una correlación menor pero significativa ($r = 0.07$, IC 95\,\% [0.01--0.12], $p = 0.014$)'' (p. 1345; traducción propia). Asimismo, no encontraron diferencias significativas entre el uso activo y el pasivo en relación con la soledad, y observaron una heterogeneidad sustancial entre los estudios. Este hallazgo es central para el presente proyecto, pues sugiere que la desconexión emocional juvenil se relaciona más con un patrón de uso compulsivo y difícil de controlar que con la cantidad de horas de exposición a las redes sociales, aunque esta última no es irrelevante.
En la misma línea, Godard y Holtzman (2024) desarrollaron un meta-análisis de 141 estudios con aproximadamente 145,000 participantes, centrado en comparar los efectos del uso activo (por ejemplo, crear contenido) frente al uso pasivo (por ejemplo, desplazarse por el contenido) de las redes sociales sobre 13 resultados de salud mental, bienestar y apoyo social. La mayoría de los tamaños de efecto resultaron insignificantes ($|r| < .10$), con excepción de las asociaciones entre el uso activo y un mayor apoyo social en línea ($r = .34$), mayor bienestar ($r = .15$), mayor afecto positivo ($r = .11$) y mayores síntomas de ansiedad ($r = .12$), y entre el uso pasivo y un mayor apoyo en línea ($r = .15$). Los análisis de moderación mostraron que el uso pasivo se asoció con peores resultados emocionales en los contextos generales de redes sociales, pero no dentro de los grupos en redes sociales. Este hallazgo complementa al anterior: los efectos de las redes sociales son, en general, pequeños y dependen del tipo de interacción y del contexto en que esta ocurre.

\section{Tipo de interacción digital como factor determinante}
Más allá de establecer que existe una relación entre uso de redes sociales y bienestar emocional, resulta necesario identificar con precisión qué características de la interacción digital explican dicha relación.
Liu et al. (2026) realizaron una revisión sistemática siguiendo las directrices PRISMA, en la cual analizaron 21 estudios empíricos publicados entre 2000 y 2024 sobre la asociación entre el uso de redes sociales y el aislamiento social en jóvenes adultos. Los autores organizaron los hallazgos en tres categorías: las motivaciones y patrones de uso de las plataformas, los efectos de los canales de comunicación textuales frente a los visuales, y la ausencia de asociación entre redes sociales y aislamiento social. La mayoría de los estudios (15) reportó efectos adversos vinculados a patrones problemáticos de uso, como el uso excesivo o adictivo, el uso pasivo y el uso compulsivo; una minoría (4) encontró asociaciones positivas relacionadas con un uso intencionado o con características específicas de las plataformas, y dos estudios no encontraron asociación significativa. Estos resultados evidencian que la relación entre redes sociales y aislamiento no es uniforme, sino que depende de la forma en que se utilizan.
En una revisión crítica, Verduyn et al. (2017) encontraron relaciones negativas entre el uso pasivo de las redes sociales y el bienestar subjetivo, y relaciones positivas entre el uso activo y el bienestar, siendo la primera más robusta que la segunda. Según los autores, el uso pasivo provoca comparaciones sociales y envidia, mientras que el uso activo, cuando predice el bienestar, parece hacerlo al generar capital social y sentimientos de conexión social. Un estudio exploratorio realizado por Tian (2025) con 150 participantes chinos de entre 18 y 57 años, si bien de alcance metodológico limitado por el tamaño y la composición de su muestra, aporta evidencia consistente con esta revisión: mediante el \textit{Social Media Activity Questionnaire} (SMAQ), encontró que el uso activo se correlacionó negativamente con la comparación social, la depresión, la ansiedad y el estrés, mientras que el uso pasivo mostró el patrón inverso.
En conjunto, esta evidencia permite identificar un primer factor relevante para el diseño de la aplicación: los mecanismos que fomentan una interacción activa, directa y recíproca entre los usuarios tienden a asociarse con mejores resultados emocionales que aquellos que promueven el consumo pasivo de contenido, aunque los tamaños de efecto son pequeños y, en el caso específico de la soledad, el meta-análisis de Feng et al. (2026) no encontró diferencias entre ambos tipos de uso.

\section{Conexión abundante versus vínculo significativo}
Un segundo eje de análisis se refiere a la naturaleza cualitativa de los vínculos que las plataformas digitales permiten construir. Desde la definición de Hawkley y Cacioppo (2010), la soledad depende no solo de la cantidad de relaciones, sino especialmente de su calidad. Bajo esta perspectiva, un joven puede acumular cientos de contactos digitales y experimentar, al mismo tiempo, un nivel elevado de soledad, en la medida en que esos contactos no satisfacen su necesidad de conexión genuina. En una revisión narrativa sobre la evolución de las redes sociales en los últimos 20 años, Smith et al. (2021) concluyen que, en lugar de preguntarse si las redes sociales generan soledad o pertenencia en los jóvenes, resulta más útil preguntarse en qué circunstancias y para quiénes ocurren estos resultados, ya que la tecnología puede tanto disminuir como fortalecer el sentido de bienestar social de los jóvenes según la naturaleza de su uso.
Esta idea encuentra respaldo empírico en el estudio longitudinal de Burke y Kraut (2016), quienes, a partir de encuestas aplicadas a 1,910 usuarios de Facebook de 91 países combinadas con los registros de su actividad en la plataforma, encontraron que la relación entre el uso de Facebook y el bienestar depende del tipo de comunicación y de la fuerza del vínculo. Recibir comunicación redactada y dirigida específicamente a la persona (\textit{composed communication}) por parte de vínculos fuertes se asoció con mejoras en el bienestar, mientras que la retroalimentación de un solo clic, como los \textit{likes}, y la lectura de publicaciones dirigidas a una audiencia amplia no mostraron dicha asociación. Estos resultados ofrecen un marco explicativo para comprender por qué las plataformas actuales, orientadas a maximizar el número de conexiones y las interacciones rápidas, tienden a producir relaciones numerosas pero superficiales.
Evidencia cualitativa reciente ilustra cómo los propios usuarios experimentan esta tensión. Parry et al. (2026) analizaron, mediante métodos computacionales y cualitativos, las discusiones de más de 26,000 miembros activos de r/NoSurf, una comunidad en línea orientada a la desconexión digital, y encontraron que sus miembros describen la tecnología simultáneamente como una vía de escape y como una fuente de angustia. Un tema recurrente fue lo que los autores denominan la ilusión de la amistad significativa en línea: los miembros describían las amistades digitales como carentes de profundidad, reciprocidad y autenticidad, por ser con frecuencia unilaterales, de bajo esfuerzo y pasivas, y señalaban que las redes sociales fomentan una ``ilusión'' de conexión en la que la persona se siente rodeada de gente, pero permanece emocional o socialmente aislada.

\section{Resultados de la encuesta diagnóstico aplicada al grupo objetivo}
\textit{Esta sección se completa una vez aplicada la encuesta de diagnóstico descrita en el apartado 1.6.2 del presente documento, dirigida a jóvenes de 16 a 28 años. Se recomienda estructurar la presentación de resultados en los siguientes puntos:}

\begin{itemize}
    \item Caracterización de la muestra (edad, género, plataformas de uso frecuente).
    \item Nivel percibido de desconexión emocional al interactuar con redes sociales actuales.
    \item Percepción sobre la calidad versus la cantidad de contactos digitales.
    \item Valoración de elementos como métricas públicas (likes, seguidores, vistas) en la experiencia emocional del usuario.
    \item Interés y expectativas frente a una herramienta centrada en micro interacciones breves y significativas.
\end{itemize}

Los resultados obtenidos deberán contrastarse con los hallazgos reportados en las secciones 4.1 a 4.3, de manera que el diagnóstico no se sustente únicamente en la literatura, sino también en evidencia primaria recogida en el contexto local del proyecto.

\section{Factores identificados que limitan la generación de vínculos significativos}
A partir de la evidencia revisada, es posible sintetizar los siguientes factores que contribuyen a la desconexión emocional de los jóvenes en su interacción con las plataformas digitales actuales:

\begin{enumerate}
    \item \textbf{Uso problemático más que uso extendido.} La soledad se asocia con mayor fuerza a un patrón de uso problemático, compulsivo y difícil de controlar ($r = 0.29$) que al tiempo total de exposición a las redes sociales ($r = 0.07$) (Feng et al., 2026).
    \item \textbf{Predominio del consumo pasivo sobre la interacción activa.} El uso pasivo de las redes sociales se ha asociado con menor bienestar subjetivo, al provocar comparaciones sociales y envidia, mientras que el uso activo tiende a asociarse con mayor bienestar y apoyo social en línea (Verduyn et al., 2017; Godard \& Holtzman, 2024; Tian, 2025). No obstante, estos efectos son pequeños y dependen del contexto, y no se han observado diferencias entre ambos tipos de uso respecto a la soledad (Feng et al., 2026).
    \item \textbf{Ausencia de mecanismos que prioricen la calidad del vínculo.} Las aplicaciones sociales populares incorporan elementos de diseño orientados a prolongar el tiempo de uso, como el desplazamiento infinito y los ``me gusta'' como recompensa social (Montag et al., 2019), y privilegian formas de interacción rápidas, como la retroalimentación de un solo clic, que no se asocian con mejoras en el bienestar; dichas mejoras se observan, en cambio, cuando la comunicación es personalizada y proviene de vínculos cercanos (Burke \& Kraut, 2016).
    \item \textbf{Presencia de métricas de comparación social.} Los elementos de validación pública propios de las redes sociales convencionales (likes, seguidores, vistas) funcionan como mecanismos de recompensa y comparación social, y este proceso de comparación puede conducir a una menor autoestima (Montag et al., 2019). La comparación social constituye, además, uno de los principales mecanismos por los que las redes sociales perjudican a sus usuarios (Clark et al., 2018; Verduyn et al., 2017).
    \item \textbf{Percepción de superficialidad en los vínculos digitales.} Los propios usuarios identifican una brecha entre la cantidad de contactos digitales disponibles y la calidad de conexión emocional que estos les proporcionan, y describen las amistades en línea como carentes de profundidad, reciprocidad y autenticidad (Parry et al., 2026), lo que resulta coherente con una definición de la soledad centrada en la calidad de las relaciones (Hawkley \& Cacioppo, 2010).
\end{enumerate}
